
\section{Method}
\label{sec:method}

\subsection{Problem Formulation}
\label{sec:problem-def}

Existing approaches formulate tool routing as direct selection from a catalog.
Given a query $q$, a model selects tools:

\begin{equation}
    g(q) \rightarrow \{t_1, \ldots, t_k\} \subseteq \mathcal{T}
\end{equation}

This requires the model to simultaneously identify relevant tools and produce a valid execution order, with no structural guarantee on the result.

We instead formulate routing over typed entities.
Let $\mathcal{E} = \{e_1, \ldots, e_n\}$ denote entity types in a domain (e.g., \texttt{Pod}, \texttt{Namespace}, \texttt{Deployment}).
Each tool is a typed transformation:

\begin{equation}
    t_k : e_i \rightarrow e_j,
    \qquad
    e_i,e_j \in \mathcal{E}.
    \label{eq:tool-transform}
\end{equation}

The objective is to recover an ordered tool sequence $[t_1,\ldots,t_l]$ that transforms a source entity into the desired target, where the output type of every tool matches the input type of the next.

\subsection{Typed Composition Graph}
\label{sec:typed-graph}

The tool registry induces a directed graph

\[
G=(\mathcal{E},\mathcal{T}),
\]

where entity types are nodes and tools are directed edges.
An edge from $e_i$ to $e_j$ exists whenever a tool $t_k:e_i\rightarrow e_j$ is available.
Multiple tools may connect the same pair, and tools with multiple input or output types generate multiple edges.

In this representation, a valid tool invocation corresponds to a graph edge and a valid multi-step workflow corresponds to a graph path.

\subsection{Typed Composition Routing}
\label{sec:tcr}

Routing is decomposed into two stages.
First, the LLM predicts source and target entity types:

\begin{equation}
    f(q)
    \rightarrow
    (\hat e_{\mathrm{src}},
     \hat e_{\mathrm{tgt}})
    \label{eq:type-prediction}
\end{equation}

Second, deterministic graph search recovers the shortest valid composition:

\begin{equation}
    \pi(q)
    =
    \operatorname{Path}
    \left(
        G,
        f_{\mathrm{src}}(q),
        f_{\mathrm{tgt}}(q)
    \right)
    \label{eq:routing-function}
\end{equation}

where $\operatorname{Path}$ returns the shortest tool sequence connecting the predicted types via breadth-first search.
If no valid path exists, the system returns an empty result rather than generating an invalid composition.

Two structural properties follow directly.

\paragraph{Proposition 1 (Structural validity).}
If $\pi(q)$ returns a non-empty sequence $[t_1, \ldots, t_l]$, then (i) every $t_i \in \mathcal{T}$ belongs to the registry, and (ii) for every adjacent pair $(t_i, t_{i+1})$, the output type of $t_i$ matches the input type of $t_{i+1}$.

\emph{Proof.}
Each edge in $G$ is a registered tool, giving (i).
Adjacent edges share a node (entity type), giving (ii) by construction. \hfill $\square$

\paragraph{Corollary (Zero hallucinated tools).}
$\pi(q)$ cannot return a tool $t \notin \mathcal{T}$, since graph search traverses only registered edges.

\subsection{Graph-Constrained Reasoning}
\label{sec:graph-constraint}

Once a target type has been predicted, reverse reachability eliminates every source type that cannot reach that target through any valid composition.
We quantify this using \emph{entity pruning}:

\begin{equation}
    \mathrm{pruning}(e_{\mathrm{tgt}})
    =
    1-
    \frac{
        |R^{-1}(e_{\mathrm{tgt}})|
    }{
        |\mathcal{E}|
    }
    \label{eq:entity-pruning}
\end{equation}

where $R^{-1}(e_{\mathrm{tgt}})$ is the set of entity types from which $e_{\mathrm{tgt}}$ is reachable.
A pruning value of 0.87 means 87\% of source candidates are structurally eliminated.
This reduction depends only on graph topology and applies equally to any type prediction model.

\subsection{Recall Decomposition}
\label{sec:recall-decomp}

Because routing decomposes into type prediction followed by graph search, end-to-end recall can be conditioned on prediction correctness:

\begin{equation}
    R
    =
    P_c R_c
    +
    (1-P_c) R_{\text{wrong}}
    \label{eq:recall-decomposition}
\end{equation}

where $P_c$ is the probability that both predicted types are correct,
$R_c$ is recall conditioned on correct predictions,
and $R_{\text{wrong}}$ is recall conditioned on incorrect predictions.

This decomposition separates LLM quality ($P_c$) from graph quality ($R_c$) and provides two independent levers for improvement: better entity type prediction improves $P_c$; richer graph connectivity improves $R_{\text{wrong}}$ (partial recovery under errors).
